\section{Quantitative confinement of the original supports}
\label{app:confinement}
This appendix proves Proposition~\ref{prop:reduction}. Its first two
lemmas turn a small lattice-moment deficit and a small smoothing loss
into explicit distances. The remaining argument applies them to the
contact equations and records the constants needed for the theorem.

\subsection{The deficit in Esseen's moment inequality}
We use a variance-weighted decomposition to quantify the deficit;
compare \citet[Lemma~A.9]{HeCheng2026} for stability of a single lattice law.
Put $c=\ce$, $p=\pe$, and $\tau=\te$ within this subsection.
Let $Y$ be centered, nondegenerate, and supported on a translate of
$h\Z$. Its oriented deficit is
\[
 \delta_Y=\frac{c\E|Y|^3-\E Y^3-3h\E Y^2}{h\E Y^2}.
\]
Write $P_-$ for the law of $-Y$ on $(0,\infty)$, $P_+$ for the
law of $Y$ there, and $m=\E Y_+=\E Y_->0$.
The centered two-point law $P_{u,v}$ at $-u,v$ has masses
$v/(u+v),u/(u+v)$. Direct integration gives
\[
 \Law(Y)=\Pp(Y=0)\delta_0+\int P_{u,v}\,\lambda(du,dv),
 \qquad\lambda(du,dv)=\frac{u+v}{m}P_-(du)P_+(dv).
\]
Its variance-weighted mixing measure
$\omega(du,dv)=uv\lambda(du,dv)/\E Y^2$ is a probability measure.
Set $n=(u+v)/h\in\{1,2,\ldots\}$ and $a=u/(u+v)$.
The identity
\[
 c\{a^2+(1-a)^2\}-(1-2a)=3+2c(a-p)^2
\]
gives the exact formula
\begin{equation}\label{eq:lattice-deficit}
 \delta_Y=\int\{3(n-1)+2cn(a-p)^2\}\,d\omega.
\end{equation}
This proves the oriented inequality and, by reflection, Esseen's
inequality \eqref{eq:esseen-moment}, including its equality case.

If $\delta_Y<3$, some component has $n=1$. Its two points are the
unique nearest lattice points straddling zero, say $-ah,(1-a)h$;
in particular this translate does not contain zero. Its variance
weight is at least $1-\delta_Y/3$. Every other component has variance
at least $h^2a(1-a)$, so this variance weight is at most its ordinary
mixture mass. Thus, for a centered law $Q$,
\begin{gather}
 \Law(Y)=(1-r_Y)\Law(h(B_a-a))+r_YQ,\quad r_Y\le\delta_Y/3,\notag\\
 |a-p|\le\sqrt{\frac{\delta_Y}{2c(1-\delta_Y/3)}},\quad
 h^2a(1-a)\le\E Y^2\le\frac{h^2a(1-a)}{1-\delta_Y/3}.   \label{eq:nearest-law}
\end{gather}
If $|Y|\le M$, coupling the mixture and then the two Bernoulli laws
bounds their Wasserstein distance by
\[
 W_1(\Law(Y),\Law(h(B_p-p)))
 \le\frac{2M}{3}\delta_Y
       +2h\sqrt{\frac{\delta_Y}{2c(1-\delta_Y/3)}}.
\]
Here $W_1$ is the infimum of $\E|Y-Z|$ over couplings.

For a collection of centered $h$-lattice variables, use one common
orientation and set
\[
 V=\sum_i\E Y_i^2,\quad B=\sum_i\E|Y_i|^3,\quad
 K=\sum_i\E Y_i^3,\quad \Delta_L=(cB-K-3hV)/(hV).
\]
Degenerate coordinates have zero variance weight. Applying the preceding
bound on $\delta_{Y_i}\le1$, the crude bound $M+h$ on its complement,
and Cauchy--Schwarz gives
\begin{equation}\label{eq:lattice-W1}
 \frac1V\sum_i\E Y_i^2 W_1(\Law(Y_i),\Law(h(B_p-p)))
 \le(5M/3+h)\Delta_L+h\sqrt{3/c}\sqrt{\Delta_L}.
\end{equation}
There are also useful moment-ratio estimates. Put
\[
 a_D=\sqrt{(1+D/3)D/(2c)},\quad
 b_D=\tau D/3+D/c+(2/c)a_D,\quad k_D=D/(3c)+2a_D.
\]
If $\Delta_L\le D$, then
\begin{equation}\label{eq:lattice-ratios}
 \left|\frac B{hV}-\tau\right|\le b_D,\qquad
 \left|\frac K{hV}-\frac1c\right|\le k_D,\qquad \frac B{hV}\ge\frac12.
\end{equation}
Indeed \eqref{eq:lattice-deficit} gives
$\E_\omega(n-1)\le D/3$,
$\E_\omega[n(a-p)^2]\le D/(2c)$, and
$|\E_\omega[n(a-p)]|\le a_D$.
Expand $n\{a^2+(1-a)^2\}$ and $n(1-2a)$ about $p$.
The last inequality follows directly from $n\ge1$.
All these formulas allow arbitrary variance weights, so they also hold
when each sum is multiplied by a common positive factor.

\subsection{Rounding at an approximate period}
For the selected array put $Y_i=X_i/\ell$.
Then $|Y_i|\le8$, $\E Y_i^2\le9$, and
\[
 \ell^2\sum_i\E Y_i^2=\ell^2\sum_i\E|Y_i|^3=1.
\]
At a nonzero frequency $u$, take the phase of
$\E e^{iuY_i}$, using any phase if the expectation vanishes.
Round $Y_i$ to the nearest point of the corresponding lattice of span
$h=2\pi/|u|$ and then center the rounded variable, calling it $Q_i$.
The inequality $|v|\le(\pi/2)|e^{iv}-1|$ for $|v|\le\pi$ gives
\begin{equation}\label{eq:rounding}
 e:=\ell^2\sum_i\E|Y_i-Q_i|^2\le h^2g(u)/8,\qquad |Q_i|\le8+h.
\end{equation}
In detail the uncentered squared rounding error is at most
$(h^2/16)\E|e^{iuY_i}-e^{i\theta_i}|^2
=(h^2/8)(1-|\E e^{iuY_i}|)$, which is no larger than the expression
using $1-|\E e^{iuY_i}|^2$. Centering decreases that error.
If $M=8+h$, Cauchy--Schwarz yields
\begin{align}
 \ell^2\sum_i|\E Q_i^2-\E Y_i^2|&\le2\sqrt e+e,\notag\\
 \ell^2\sum_i|\E|Q_i|^3-\E|Y_i|^3|,\quad
 \ell^2\sum_i|\E Q_i^3-\E Y_i^3|
                &\le(3M/2)\sqrt e(2+\sqrt e).           \label{eq:rounding-moments}
\end{align}
For the cubic bounds use
$||x|^3-|y|^3|,|x^3-y^3|\le(3M/2)|x-y|(|x|+|y|)$.
These bounds depend on total variance, not on the number of coordinates.

\subsection{Comparing two small smoothing losses}
\begin{lemma}\label{lem:support-spacing}
Let $F$ be the CDF of a finite positive measure and let
$J=F*\operatorname{Unif}[-h/2,h/2]$, $h>0$.
Suppose $J(x)$ differs from an affine function of slope $q>0$ by at
most $\varepsilon$ on $[-D_0-2h,3h]$.
Put $H(a)=J(a+h/2)-F(a)$. If $0\le d\le D_0$ and
$H(0),H(-d)\le s$, then, for $0<r_0\le h/4$ and
$K_0=\lceil D_0/h\rceil$,
\begin{equation}\label{eq:spacing-bound}
 \dist(d,h\Z)\le3r_0+
 \frac{4s+(5+4K_0)\varepsilon+4K_0h\varepsilon/r_0}{q}.
\end{equation}
\end{lemma}
\begin{proof}
For the measure $\nu$ of $F$, write
\[
 w_a(t)=(1-(t-a)/h)\ind_{(a,a+h)}(t),\qquad
 H(a)=\int w_a(t)\nu(dt)\ge0.
\]
For $0<d<h$, define the nonnegative trapezoid
\[
 T_d(t)=\frac1h\int_{h/2-d}^{h/2-d/2}
                  \ind_{\{|x-t|<h/2\}}\,dx.
\]
Its integral against $\nu$ equals $J(h/2-d/2)-J(h/2-d)$.
Splitting at $-d,-d/2,0,h-d,h-d/2,h$ gives the pointwise bound
\[
 \min\{1,2(h-d)/d\}\,T_d(t)\le w_0(t)+w_{-d}(t).
\]
Integrating and applying the affine approximation to the two $J$ values gives
\begin{equation}\label{eq:two-shifts}
       q\min(d/2,h-d)\le H(0)+H(-d)+2\varepsilon.
\end{equation}
The pointwise bound also holds at the endpoints, with the open intervals
in $w_a$, so atoms require no continuity assumption on $F$.

To transfer over several periods, let
$\overline H_{r_0}(a)=r_0^{-1}\int_a^{a+r_0}H(t)dt$.
Integrating the definition of the uniform convolution gives
\begin{align*}
 \overline H_{r_0}(a+h)-\overline H_{r_0}(a)
 ={}&r_0^{-1}\int_a^{a+r_0}[J(t+3h/2)-J(t+h/2)]dt\\
 &-(h/r_0)[J(a+r_0+h/2)-J(a+h/2)].
\end{align*}
The affine terms cancel, so its magnitude is at most
$2\varepsilon(1+h/r_0)$.
Moreover
$H(a+t)\le H(a)+(t/h)\nu((a+t,a+t+h))$ and the mass is at most
$2(qh+\varepsilon)$, by the affine bound on $J$. Hence
$\overline H_{r_0}(a)\le H(a)+qr_0+\varepsilon r_0/h$.
Write $d=kh+d_0$, $0\le d_0<h$. Transfer from $-d$ to $-d_0$
using the preceding difference bound. If $d_0$ is within $2r_0$
of an endpoint, the conclusion is immediate. Otherwise choose a point
in the averaging interval where $H$ is at most its average and apply
\eqref{eq:two-shifts}, using $H(0)\le s$.
The threshold displacement is at most $r_0$; the resulting bound is
\eqref{eq:spacing-bound}. All arguments of $J$ stay in the stated window.
\end{proof}

\subsection{The scale and the maximizing threshold}
We prove Proposition~\ref{prop:reduction}. Retain the notation of
Section~\ref{sec:selection} for the actual selected maximizer.
For scaled coordinates write
\[
 Y_i=X_i/\ell,\quad v_i=\E Y_i^2,\quad M_i=\E[Y_i|Y_i|],
 \quad\theta=\ell^{-1}\sum_i\E X_i^3,\quad |\theta|\le1.
\]
Let $\rho,r,T,\eps,h$ be as in Proposition~\ref{prop:jitter}.
We give the error quantities explicitly so their later specialization
requires no unspecified constants:
\begin{equation}\label{eq:confinement-errors}
\begin{gathered}
\begin{aligned}
 u&=\rho+\ell/2,& D&=16u+256r,\\
 H_h&=72r+5b_D,& Z^2&=10(u+18r),\\
 S_0&=2\rho+18r+460\ell^2,&
 A_0&=2\rho+24\{4\ell^2+(Z+32\ell)^2/2+\ell/5\},\\
 \gamma&=300\sqrt{A_0}+11S_0+131A_0,&
 W_0&=21D+3\sqrt D+42r,
\end{aligned}\\
\begin{aligned}
 E_P&=14000(u+18r+\eta)+1100Z^2+4000\ell+4000\ell^2+1100\eta,\\
 E_{\rm end}&=E_P+6100\gamma+700\gamma^2+500H_h.
\end{aligned}
\end{gathered}
\end{equation}
Here $a_D,b_D,k_D$ are those in \eqref{eq:lattice-ratios}.
The proof below uses only
\begin{equation}\label{eq:confinement-gates}
 \begin{gathered}
 r\le10^{-6},\quad D\le10^{-3},\quad A_0\le9/64,\quad
 \gamma\le1/100,\\
 E_P\le10^{-4},\quad E_{\rm end}\le1/2000,\quad W_0\le1/4.
 \end{gathered}
\end{equation}
These will follow from \eqref{eq:reduction-parameters} at the end.

The case $h=0$ would give $R\le\phz/6+\rho<\be$ by
\eqref{eq:jitter}, so $h>0$. Initially $h\le128\pi<404$ and
$g(2\pi/h)\le r^2/4$. Equations
\eqref{eq:rounding}--\eqref{eq:rounding-moments} give
$\sqrt e\le101r$, $|Q_i|\le412$.
If $V_Q=\ell^2\sum\E Q_i^2$ and
$B_Q=\ell^2\sum\E|Q_i|^3$, then
$V_Q\ge1-202r$ and $B_Q\le1+125000r$.
Since $B_Q/(hV_Q)\ge1/2$, it follows that
\[
 h\le2(1+125000r)/(1-202r)<5/2.
\]
Apply the same rounding bounds with this improved span. Writing
$K_Q=\ell^2\sum\E Q_i^3$, we obtain
\begin{equation}\label{eq:rounded-data}
 \sqrt e\le r,\quad|V_Q-1|\le3r,\quad
 |B_Q-1|\le32r,\quad|K_Q-\theta|\le32r.
\end{equation}
The lattice inequality in either orientation implies
\begin{equation}\label{eq:span-envelope}
                3h+|\theta|\le c+256r.
\end{equation}
The moment-error allowance is $[32(c+1)+9h]r<256r$.

Unsmoothing at the actual threshold $z$, with shift $h\ell/2$, gives
\begin{equation}\label{eq:global-envelope}
 R\le\phi(z)\{h/2+\theta(1-z^2)/6\}+u.
\end{equation}
The Gaussian increment has error at most $\phz h^2\ell^2/8$,
and the skew correction adds at most $h\ell^2/18$;
together they are below $\ell^2/2$.
Since $\phi(z)$ and $|(1-z^2)\phi(z)|$ are at most $\phz$,
\begin{equation}\label{eq:near-max-envelope}
 3h+|\theta|>c-16u,\qquad R\le\be+u+18r.
\end{equation}
It follows that $h>3/2$. Orient the rounded variables by
$\operatorname{sign}\theta$, arbitrarily if $\theta=0$.
Equations \eqref{eq:rounded-data}--\eqref{eq:near-max-envelope}
bound their aggregate lattice deficit by $D$, its denominator
$hV_Q$ being greater than one. The moment-ratio estimates yield
\begin{equation}\label{eq:near-esseen}
 |h-\he|\le H_h,\qquad
 \bigl||\theta|-1/\sqrt{10}\bigr|\le34r+3k_D+H_h/6.
\end{equation}
For $D\le10^{-3}$,
$a_D\le\sqrt D/3$,
$b_D\le D/2+\sqrt D/9$, and
$k_D\le D/18+2\sqrt D/3$.
Thus $H_h<1/40$ and the second error is below $1/10$.
Since $31/16<\he<41/21$, this proves
\begin{equation}\label{eq:span-range}
                  19/10<h<2,\qquad|\theta|>1/5.
\end{equation}
If $\theta<0$, the envelope in \eqref{eq:global-envelope} has its
maximum at zero because $h>|\theta|$; it would imply
$|\theta|<128r+8u=D/2$, a contradiction. Hence $\theta>0$.
Finally
$\be<(\be+18r)e^{-z^2/2}+u$, so
\begin{equation}\label{eq:threshold-bound}
           |z|\le Z,\qquad0\le K-\be\le u+18r+\eta.
\end{equation}
For this last bound use $-\log(1-x)\le2x$ for $x\le1/2$.

\subsection{Deletion with the same period and the support spacing}
Deleting coordinate $i$ subtracts at most $\ell^2$ from $g$.
Thus its high-frequency product bound is at most $e^{1/2}<2$ times
the full-row bound. Keep the original $h$ for every deletion.
Write $\sigma_i^2=1-s_i\ge3/4$, $\kappa_- =\sum_{j\ne i}\E X_j^3$,
let $J_i$ be the CDF of that deletion plus the same jitter, and put
\[
 G_i^*(x)=\Phi(x/\sigma_i)
       +\frac{\kappa_-}{6\sigma_i^3}(1-(x/\sigma_i)^2)\phi(x/\sigma_i).
\]
The logarithmic remainder in Appendix~\ref{app:fourier} remains
$(49/3)\ell^2t^4$. The low-frequency error is bounded by
$\ell^2e^{-t^2/5}\{(49/3)t^4+t^6/72\}$, whose positive integral
is below $206\ell^2$. The jitter-factor error is below $\ell^2$.
For $a=\delta/\ell$, the Gaussian comparison tail is at most
$4/(3a^2)+\ell/(2a)<5500\ell^2$; use
$\int_a^\infty t^2e^{-3t^2/8}dt\le3/a$.
Also $G_i^*$ is Lipschitz with constant below $1/2$.
The weaker signed smoothing bound therefore gives
\begin{equation}\label{eq:deleted-jitter}
 \|J_i-G_i^*\|_\infty
 \le6000\ell^2+24r\ell/\delta^2
       +2\log(T/\delta)e^{-\eps^2/(2\ell^2)}+20\ell/T
 \le2\rho\ell.
\end{equation}
The last inequality uses the four allocations in
\eqref{eq:jitter-parameters}.

At any original support point $y=\ell t$, the contact equality and
\eqref{eq:gaussian-global} give
\begin{equation}\label{eq:contact-flat}
 |G_i(z-y)-\Phi((z-y)/\sigma_i)-R\ell|\le450\ell^3.
\end{equation}
For the constant, $|t|\le8$, $v_i,|M_i|\le9$,
$\E|Y_i|^3\le27$, and $K\le23/50$.
The Gaussian remainder contributes less than $51\ell^3$,
the quadratic $K$ term less than $51\ell^3$, and the remaining
cubic, moment, and linear terms less than $348\ell^3$.
Equations \eqref{eq:span-envelope}, \eqref{eq:deleted-jitter}, and
\eqref{eq:contact-flat} then imply
\begin{equation}\label{eq:smoothing-loss}
 0\le[J_i(z-y+h\ell/2)-G_i(z-y)]/\ell\le S_0.
\end{equation}
The scale corrections cost at most $10\ell^2$ after division by
$\ell$: use $\sigma_i^{-1}-1\le s_i$,
$\sigma_i^{-3}-1\le3s_i$, and
$|\kappa_-/\ell|\le\theta+27\ell^2$.

Fix a support point $x$. The function
$\ell^{-1}J_i(z-x+\ell a)$ differs from an affine function of slope
$\phz$ by at most $A_0$ on $[-16-2h,3h]$.
Indeed $|a|\le24$, $|z-x+\ell a|\le Z+32\ell$, and
\[
                 |(G_i^*)'(w)-\phz|\le4\ell^2+w^2/2+\ell/5.
\]
Integrate and use \eqref{eq:deleted-jitter}.
Apply Lemma~\ref{lem:support-spacing} to the scaled positive measure,
using \eqref{eq:smoothing-loss}, $q=\phz$, $K_0\le11$, and
$r_0=\sqrt{A_0}\le h/4$. Its bound is at most
$3\sqrt{A_0}+(8/3)(4S_0+49A_0+110\sqrt{A_0})<\gamma$.
Thus for every pair of support points of each original coordinate,
\begin{equation}\label{eq:arithmetic-spacing}
                   \dist((y-x)/\ell,h\Z)\le\gamma.
\end{equation}

\subsection{The support polynomial and the two classes}
Define
\[
 P_i(t)=t^3-c|t|^3+(3c/2)t^2+3(cM_i-v_i)t.
\]
Every point of the full scaled support satisfies
\begin{equation}\label{eq:support-polynomial}
                             P_i(t)\ge-E_P.
\end{equation}
To verify this with constants, divide \eqref{eq:deletion} by
$\phz\ell^3/6$ and use \eqref{eq:gaussian-local}.
Its Taylor remainder divided by $\ell^3$ is below $224\ell$.
Furthermore $|\phi''(z)+\phz|\le(3/2)\phz Z^2$,
since $\|\phi''''\|_\infty\le3\phz$ and $\phi'''(0)=0$.
The respective error bounds are
\[
 13184(K-\be),\quad1092Z^2,\quad3584\ell,\quad
                         3500\ell^2,\quad1008\eta.
\]
For the first, use
$|-|t|^3+(3/2)t^2+3M_it|\le824$ and $6/\phz<16$.
The last two come from $6R\ell s_i^2+\eta b_i+4\eta\ell s_i$.
Equation~\eqref{eq:threshold-bound} makes their sum at most $E_P$.

For a nondegenerate coordinate let $-a,b$ be its support extremes.
Centering gives $v_i\le ab$ and $|M_i|\le v_i$.
If $a,b\ge1/100$, divide the endpoint inequalities by $a,b$ and
add. The linear terms cancel, giving
\[
 (c+1)a^2+(c-1)b^2\le(3c/2)(a+b)+E_P(a+b)/(ab).
\]
The left side is at least $3(a+b)^2$, because
$(c^2-1)/(2c)=3$. Therefore
$a+b\le c/2+E_P/(3ab)\le25/8+1/3<7/2$.
If $a<1/100$ and $a+b\ge7/2$, then $b>3.49$ and
\[
 P_i(b)\le b^2[-(c-1)b+3c/2+3(c-1)a]<-1,
\]
contradicting \eqref{eq:support-polynomial}. If $b<1/100$, use
$P_i(-a)\le a^2[-(c+1)a+3c/2+3(c+1)b]<-1$.
Every support diameter is consequently below $7/2$.

Partition coordinates according as their scaled support diameter is
less than $h/2$ or at least $h/2$. In the first class,
\eqref{eq:arithmetic-spacing} implies diameter at most $\gamma$,
so centering gives $|Y_i|\le\gamma$.
In the second class the diameter differs from $h$ by at most $\gamma$:
zero and two multiples are excluded by \eqref{eq:span-range} and
the diameter bound. Every support point lies within $2\gamma$ of
an extreme. Split at the midpoint, use its upper-cluster indicator
$B_i$, and take conditional means. This gives exactly
\begin{equation}\label{eq:cluster-representation}
 Y_i=h_i(B_i-p_i)+U_i,\quad\E(U_i\mid B_i)=0,\quad
 |h_i-h|\le5\gamma,\quad |U_i|\le2\gamma.
\end{equation}
Both cluster probabilities are positive, because both extremes are
support points and the clusters are separated.

We next bound those probabilities away from zero and one.
On $[-8,8]$, $|P_i'|<1800$. Moving an extreme to its cluster mean
costs at most $3600\gamma$. Conditional centering gives
$|M_i-M_i^0|\le\E U_i^2\le4\gamma^2$ and
$v_i-v_i^0=\E U_i^2\le4\gamma^2$, where the superscript denotes
the pure two-point law. The first bound uses the Lipschitz derivative
of $x|x|$ and also holds if a cluster crosses zero.
Their polynomial contribution is below $700\gamma^2$.
At a pure-law endpoint the polynomial is
$h_i^3C_3(p_i)+h_i^2C_2(p_i)$, with $|C_3|<29$, $|C_2|<75/8$.
Its gap derivative for gaps at most $9/4$ is below $500$.
Thus moving to the ideal gap $\he$ leaves both endpoint values at
least $-E_{\rm end}$.

Put $t_0=3\sqrt{10}/2$. At its upper and lower endpoints those
ideal-gap values are
\[
 \he^3(1-p)^2A(p),\qquad\he^3p^2B(p),
\]
where
\begin{align*}
 A(p)&=t_0-(c-1)+4(c-1)p-6cp^2,\\
 B(p)&=t_0-3(c-1)+(8c-4)p-6cp^2.
\end{align*}
For $p\le1/4$, $B$ is increasing and
$B(p)\le(\sqrt{10}-17)/8<-1$.
If $p\in[1/100,1/4]$, the lower endpoint is below
$-6/10000<-1/2000$, since $\he^3>6$.
For $p\le1/100$, use $A(0)=-\pe<-2/5$, $A'\le21$, so
$A(p)<-19/100$. For $p\ge3/4$, $A$ is decreasing and
$A(p)\le(\sqrt{10}-49)/8<-5$, excluding $[3/4,99/100]$.
For $p\ge99/100$, $B(1)<-2$ and $|B'|\le46$ give $B(p)<-1$.
In view of $E_{\rm end}\le1/2000$, all these cases are impossible.
Hence every two-cluster probability lies in $[1/4,3/4]$.

\subsection{The extra coordinates and the parameter check}
Lattice stability and the rounding coupling give
\begin{equation}\label{eq:original-W1}
 \sum_i\ell^2v_i W_1(\Law(Y_i),\Law(h(B_{\pe}-\pe)))\le W_0.
\end{equation}
For the rounded variables \eqref{eq:lattice-W1} is at most
$(1+3r)(20D+2\sqrt D)\le21D+3\sqrt D$.
The change of variance weights has total absolute mass at most $3r$,
and the distances are at most $13$, costing $39r$.
The original variance-weighted rounding distance is at most $3r$,
by $v_i\le9$ and the squared error bound $e\le r^2$.
This proves \eqref{eq:original-W1} with the $42r$ allowance in $W_0$.

The target law has absolute first moment
$2h\pe\qe>9/10$. On the small-diameter class $|Y_i|\le\gamma$.
Its total original variance $V_s$ is therefore at most
$W_0/(9/10-\gamma)\le2W_0$.
Each remaining coordinate has scaled variance between $1/2$ and $2$,
by \eqref{eq:span-range}, \eqref{eq:cluster-representation}, and
$p_i\in[1/4,3/4]$. Their number $m$ satisfies
\begin{equation}\label{eq:core-count}
 \frac{1-2W_0}{2\ell^2}\le m\le\frac2{\ell^2},
                         \qquad m\ge1/(4\ell^2).
\end{equation}
Calling the small scaled coordinates $Z_j$,
\[
 \frac{\sum_j\Var Z_j}{mh^2}
       \le\frac{4W_0}{h^2(1-2W_0)}<3W_0.
\]
The relative gap errors in \eqref{eq:cluster-representation} are
below $3\gamma$, the residuals satisfy $|U_i|/h<2\gamma$, and
$|Z_j|/h<\gamma$.
Thus the conclusion follows once $3\gamma,3W_0\le\lambda$ and
$\ell\le1/(2\sqrt N)$.

We finish by checking these conditions and
\eqref{eq:confinement-gates}. From \eqref{eq:jitter-parameters},
$\ell\le\rho/48000$ and $r=\rho/196608$. For $\rho\le10^{-6}$,
\begin{gather*}
 D\le17\rho,\quad Z^2\le20\rho,\quad
 H_h\le43\rho+3\sqrt\rho,\quad S_0\le3\rho,\quad A_0\le500\rho,\\
 \gamma\le7000\sqrt\rho,\quad W_0\le14\sqrt\rho,
             \quad E_P\le51000\rho+15100\eta.
\end{gather*}
For $\rho\le10^{-12}$ these imply
\[
 E_{\rm end}\le5\cdot10^7\sqrt\rho+2\cdot10^4\eta.
\]
Before enlargement its $\rho$ and $\sqrt\rho$ coefficients are
at most $34300072500$ and $42701500$, which proves this last bound.
Substitute $\rho=10^{-26}\lambda^2$, $\eta=10^{-10}\lambda$:
\[
 E_{\rm end}\le7\cdot10^{-6}\lambda<1/2000,\quad
 \gamma\le7\cdot10^{-10}\lambda,\quad
 W_0\le1.4\cdot10^{-12}\lambda.
\]
All the stated conditions follow, including $E_P\le10^{-4}$.
The representations and the variance estimates concern the original
coordinates throughout, proving Proposition~\ref{prop:reduction}.
